% !TEX program = xelatex
\documentclass[10pt,a4paper]{article}

% ---- Пакеты для XeLaTeX и кириллицы ----
\usepackage{fontspec}
\usepackage{polyglossia}
\setmainlanguage{russian}
\setotherlanguage{english}

% ---- Шрифт с кириллицей ----
\setmainfont{Liberation Serif}
\setsansfont{Liberation Sans}
\setmonofont{Liberation Mono}

% ---- Поля и геометрия ----
\usepackage[top=1.5cm, bottom=1.5cm, left=1.5cm, right=1.5cm]{geometry}
\usepackage{parskip}
\setlength{\parindent}{0pt}

% ---- Продвинутые ссылки ----
\usepackage{hyperref}
\hypersetup{
    colorlinks=true,
    linkcolor=blue,
    urlcolor=blue,
    citecolor=blue
}

% ---- Полезные пакеты ----
\usepackage{enumitem}
\setlist{nosep, left=0pt}
\usepackage{tabularx}
\usepackage{array}
\usepackage{booktabs}
\usepackage{xcolor}
\usepackage{titlesec}

% ---- Настройка заголовков разделов ----
\titleformat{\section}{\Large\bfseries\uppercase}{}{0em}{}[\vspace{-0.3em}\rule{\textwidth}{0.5pt}]
\titlespacing*{\section}{0pt}{0.8em}{0.3em}

\titleformat{\subsection}{\large\bfseries}{}{0em}{}
\titlespacing*{\subsection}{0pt}{0.4em}{0.2em}

% ---- Пользовательские команды ----
\newcommand{\projectEntry}[3]{\textbf{#1} \hfill \textit{#2}\par #3\par\vspace{0.3em}}

% ---- Общие данные (имя, телефон, email и т.д.) ----
% ============================================================
% ОБЩИЕ ПЕРСОНАЛЬНЫЕ ДАННЫЕ
% ============================================================
% Меняешь здесь — меняется во всех резюме

\newcommand{\PersonName}{Кузык Иван Романович}
\newcommand{\PersonPhone}{+7 (916) 574-52-99}
\newcommand{\PersonEmail}{vanya.kuzyk1@gmail.com}
\newcommand{\PersonGitHub}{github.com/vanya890}
\newcommand{\PersonCity}{г.~Москва, м.~Бульвар Дмитрия Донского}
\newcommand{\PersonCitizenship}{РФ}
\newcommand{\PersonDrivingLicense}{Категория B}

% ============================================================
% КОМАНДА ДЛЯ ЗАГОЛОВКА РЕЗЮМЕ
% ============================================================
% Использование: \makeHeader{Должность}
\newcommand{\makeHeader}[1]{%
\begin{center}
    {\Huge\bfseries \PersonName}\\[0.3em]
    {\large #1}\\[0.5em]
    \begin{tabular}{ll}
        Телефон: & \PersonPhone \\
        Email:   & \href{mailto:\PersonEmail}{\PersonEmail} \\
        GitHub:  & \href{https://\PersonGitHub}{\PersonGitHub} \\
        Город:   & \PersonCity \\
        Гражданство: & \PersonCitizenship
    \end{tabular}
\end{center}
}

% ============================================================
% ОБЩЕЕ ОБРАЗОВАНИЕ
% ============================================================
% Использование: \makeEducation
\newcommand{\makeEducation}{%
\section{Образование}

\projectEntry{Московский политехнический университет --- Факультет информационных технологий}{2022--2026}
{Бакалавриат, направление «Информатика и вычислительная техника», специализация «Киберфизические системы». ВКР «Разработка образовательного 3D-симулятора робототехники» --- отлично. Средний балл --- 5.0.}

\projectEntry{ГБПОУ «Образовательный комплекс «Юго-Запад»}{2018--2022}
{Среднее специальное образование, специальность «Информационные системы и программирование». Диплом с отличием. Демонстрационный экзамен по компетенции «Разработка компьютерных игр и мультимедийных приложений» --- отлично.}
}

% ============================================================
% ОБЩАЯ ДОПОЛНИТЕЛЬНАЯ ИНФОРМАЦИЯ
% ============================================================
% Использование: \makeFooter{Ключевые компетенции}{Личные качества}
\newcommand{\makeFooter}[2]{%
\section{Дополнительная информация}

\begin{itemize}
    \item \textbf{GitHub:} \href{https://\PersonGitHub}{\PersonGitHub} --- #1
    \item \textbf{Ключевые компетенции:} #2
    \item \textbf{Личные качества:} Обучаемость, ответственность, умение работать с большим объёмом информации, системное мышление, ориентация на результат
    \item \textbf{Водительские права:} \PersonDrivingLicense
\end{itemize}
}

\begin{document}

% ============================================================
% ЗАГОЛОВОК
% ============================================================
\makeHeader{Техник-программист / Разработчик / Специалист широкого профиля}

% ============================================================
% О СЕБЕ
% ============================================================
\section{О себе}

Разносторонний специалист с техническим образованием в области информационных систем и программирования. Имею опыт работы с широким спектром технологий: от веб-разработки и программирования до робототехники, 3D-моделирования и AI-систем. Выполнял проекты в рамках учёбы, фриланса и личных экспериментов. Быстро осваиваю новые инструменты, умею решать задачи любой сложности на стыке разных дисциплин. Ориентирован на практический результат.

% ============================================================
% ОПЫТ РАБОТЫ (ПРОЕКТНАЯ ДЕЯТЕЛЬНОСТЬ)
% ============================================================
\section{Опыт работы}

\subsection{Веб-разработка и программирование}

\projectEntry{Fullstack-разработчик (EduRoboSim — ВКР)}{2025–2026}
{Разработал образовательный 3D-симулятор робототехники на стеке \textbf{Tauri (Rust + WebView2) + React 18 + TypeScript + Three.js (R3F)}. Спроектировал собственный язык программирования (EduScript) с лексером, парсером (AST) и виртуальной машиной. Внедрил физический движок 2D OBB-коллизий, систему датчиков и редактор уровней.}

\projectEntry{Fullstack-разработчик (AI Chat Game)}{2024–2025}
{Создал веб-приложение для чата с AI-персонажами на \textbf{React 18 + TypeScript + Vite}. Реализовал систему памяти, визуальный редактор персонажей, дашборд аналитики токенов. Интегрировал поддержку AI-провайдеров: OpenAI, Ollama, LM Studio, Groq.}

\projectEntry{Разработчик (laravelShop — диплом колледжа)}{2022}
{Разработал интернет-магазин на \textbf{Laravel (PHP)} с каталогом товаров, корзиной, оформлением заказов и административной панелью. Использовал MySQL, Blade-шаблоны, Bootstrap.}

\projectEntry{SEO-оптимизация и работа с API (Stroy-Materiali-24.ru)}{2024}
{Разработал инструменты для SEO-оптимизации сайта строительных материалов. Работал с \textbf{Яндекс.Вебмастер}, \textbf{Яндекс API}, проводил анализ поисковых запросов, парсинг данных, генерацию контента с использованием AI-промптов.}

\subsection{AI, нейросети и машинное обучение}

\projectEntry{Разработка AI-агента для автоматической публикации статей}{2025}
{Создал бота для сайта, который автоматически генерировал и публиковал статьи на актуальные темы с использованием языковых моделей. Интегрировал \textbf{Claude, Yandex AI, Google AI Studio, Groq, Qwen Coder} для генерации контента.}

\projectEntry{Исследователь (AI/ML — нейросети и обучение)}{2024–2026}
{Выполнил цикл лабораторных и учебных проектов по искусственному интеллекту и машинному обучению:
\begin{itemize}
    \item Реализовал самописный \textbf{перцептрон} для распознавания чисел на Python
    \item Разработал систему определения входящих/выходящих людей из автобуса на Python с использованием \textbf{OpenCV}
    \item Реализовал управление беспилотным 2D-роботом в гоночной симуляции (обучение с подкреплением)
    \item Работал с \textbf{NumPy, Matplotlib, Scikit-learn} для численных расчётов и визуализации
\end{itemize}}

\projectEntry{Работа с моделями генерации изображений}{2024–2025}
{Настраивал пайплайны генерации изображений в \textbf{Stable Diffusion} и \textbf{ComfyUI}. Обучал собственные \textbf{LoRA-модели} на подготовленных датасетах. Работал с ControlNet, дообучением моделей.}

\projectEntry{Развёртывание локальных LLM}{2024–2025}
{Устанавливал и настраивал локальные языковые модели с использованием \textbf{Unsloth, LM Studio, Ollama, llama.cpp}. Настраивал инференс, работал с открытыми весами моделей.}

\subsection{Инженерные и CAD-системы}

\projectEntry{Работа с Компас-3D и AutoCAD}{2023–2025}
{Создавал 3D-модели и чертежи в \textbf{Компас-3D} в рамках учебных проектов. Работал с \textbf{AutoCAD} для 2D-моделирования и оформления конструкторской документации.}

\projectEntry{Моделирование микросхем и дорожек}{2023–2024}
{Выполнял разводку печатных плат и работал с принципиальными схемами в рамках учебных проектов.}

\subsection{Робототехника и техническое зрение}

\projectEntry{Лабораторные работы с ROS (Clover \& TurtleBot3)}{2024–2025}
{Выполнил цикл лабораторных работ в среде \textbf{ROS (Robot Operating System)}:
\begin{itemize}
    \item Разработал полётное задание для дрона Clover с использованием маркеров ArUco
    \item Работал с \textbf{лидарами}, дальномерами, датчиками линии
    \item Запускал симуляцию TurtleBot3 в Gazebo, работал с RViz
\end{itemize}}

\projectEntry{Система визуальной локализации робота (OpenCV)}{2024}
{Разработал desktop-приложение на \textbf{C\# Windows Forms} с использованием \textbf{OpenCvSharp} для определения положения и траектории робота по стереопотоку с двух камер.}

\projectEntry{Сборка робототехники}{2023–2024}
{Работал с \textbf{сервоприводами}, датчиками линии, лидарами, дальномерами. Собирал и программировал простых роботов в рамках учебных проектов.}

\projectEntry{Симуляции промышленных ячеек}{2024}
{Работал с симуляторами \textbf{роботов-манипуляторов}, программировал траектории движения в рамках учебных проектов.}

\projectEntry{Физические симуляции (LabVIEW)}{2023–2024}
{Создавал виртуальные приборы и моделировал физические процессы в среде \textbf{LabVIEW}.}

\subsection{VR/AR}

\projectEntry{Разработка AR-порталов}{2024–2025}
{Создавал AR-порталы на базе маркеров и отслеживания движений. Работал с \textbf{веб-AR} через камеру браузера (JavaScript/Three.js/AR.js).}

\projectEntry{Системы виртуальной реальности}{2024}
{Работал с VR-очками, настраивал системы синхронного воспроизведения на нескольких VR-устройствах.}

\subsection{1С, учётные системы и администрирование}

\projectEntry{Оператор ПК (ГБПОУ «Юго-Запад»)}{Ноябрь 2020 – Январь 2021}
{Ввод и редактирование данных в программе \textbf{1С:Библиотека}, обучение сотрудников работе с системой на базовом уровне.}

\projectEntry{Работа с 1С:Торговая точка и МойСклад}{2023–2024}
{Настраивал кассовые места в \textbf{1С:Торговая точка}, работал с номенклатурой. Имею опыт работы с сервисом \textbf{МойСклад} для ведения складского учёта.}

\subsection{Сетевая инфраструктура}

\projectEntry{Перенос серверного оборудования (Гос. учреждение)}{2024}
{Участвовал в переносе серверного оборудования, занимался коммутацией и базовым администрированием сетевой инфраструктуры.}

\subsection{Практика}

\projectEntry{Техник-программист (ГБПОУ «Юго-Запад»)}{Апрель–Июнь 2022}
{Проверка работоспособности ПК, обновление и установка ОС и ПО, настройка сетевых подключений, установка принтеров и периферийного оборудования, помощь в организации ЕГЭ.}

\projectEntry{Оператор БД (Издательство «Просвещение»)}{Май–Июль 2021}
{Заполнение базы данных информацией, описание изображений для поиска по ключевым словам, подготовка данных для редакций.}

% ============================================================
% ДОПОЛНИТЕЛЬНЫЕ ПРОЕКТЫ И ЗАДАЧИ
% ============================================================
\section{Дополнительные проекты и задачи}

\begin{itemize}
    \item \textbf{Telegram-боты:} Создавал ботов на Python (aiogram) для автоматизации задач, интеграции с API
    \item \textbf{Discord-боты:} Разрабатывал простых ботов на discord.py для модерации и команд
    \item \textbf{Парсинг сайтов:} Собирал данные с сайтов с помощью BeautifulSoup/Scrapy (Python)
    \item \textbf{RSS-фиды:} Работал с RSS-лентами для агрегации и получения статей
    \item \textbf{Автотесты:} Писал unit-тесты для проектов на Vite/Node.js (Vitest/Jest)
    \item \textbf{Оптимизация систем:} Анализировал загрузку системы, искал узкие места и паразитные процессы
    \item \textbf{Промпт-инжиниринг:} Разрабатывал структурированные промпты, JSON schema для управления выводом LLM
    \item \textbf{Обработка видео:} Базовый монтаж и цветокоррекция в \textbf{DaVinci Resolve}
    \item \textbf{Обработка фото:} Редактирование изображений с помощью AI-инструментов, оживление статичных фото
    \item \textbf{Социальные сети:} Ведение сообществ, создание контента, базовое продвижение
    \item \textbf{Написание работ на заказ:} Курсовые, диссертации и исследовательские работы по техническим направлениям
    \item \textbf{LaTeX:} Вёрстка документов, работа с формулами, таблицами, библиографией
\end{itemize}

% ============================================================
% ОБРАЗОВАНИЕ
% ============================================================
\makeEducation

% ============================================================
% НАВЫКИ
% ============================================================
\section{Навыки}

\begin{tabular}{p{3.5cm} p{12cm}}
    \textbf{Языки программирования:} & TypeScript, JavaScript, Python, PHP, C\#, SQL, Rust (базовый) \\
    \textbf{Frontend:} & React, TypeScript, Vite, Tailwind CSS, Three.js / R3F, HTML5, CSS3 \\
    \textbf{Backend:} & Node.js, Express.js, Laravel, REST API, WebSocket \\
    \textbf{Базы данных:} & MySQL, SQLite, IndexedDB, Oracle (базовый), 1С \\
    \textbf{AI/ML:} & OpenAI API, Claude, Yandex AI, Google AI Studio, Groq, Qwen Coder, Ollama, LM Studio, Unsloth, llama.cpp \\
    \textbf{Генерация изображений:} & Stable Diffusion, ComfyUI, LoRA, ControlNet \\
    \textbf{Нейросети:} & Перцептроны, CNN, OpenCV, NumPy, Scikit-learn, PyTorch (базовый) \\
    \textbf{CAD/Чертежи:} & Компас-3D, AutoCAD \\
    \textbf{Робототехника:} & ROS, Gazebo, RViz, Clover, TurtleBot3, сервоприводы, лидары, дальномеры \\
    \textbf{Симуляции:} & LabVIEW, симуляторы промышленных ячеек \\
    \textbf{VR/AR:} & Three.js/AR.js, веб-AR, VR-очки \\
    \textbf{Инструменты:} & Git, Docker, Docker Compose, Nginx, Linux \\
    \textbf{Медиа:} & DaVinci Resolve, AI-обработка фото, соцсети \\
    \textbf{Прочее:} & Парсинг (BeautifulSoup/Scrapy), Telegram/Discord боты, RSS, автотесты (Vitest/Jest), LaTeX, SEO, Яндекс.Вебмастер, Яндекс API, МойСклад, 1С:Торговая точка \\
    \textbf{Языки:} & Русский — родной, Английский — A2 (элементарный)
\end{tabular}

% ============================================================
% ДОПОЛНИТЕЛЬНАЯ ИНФОРМАЦИЯ
% ============================================================
\makeFooter{открытые проекты по веб-разработке, робототехнике, AI/ML, симуляциям}{Разработка ПО, веб-программирование, AI/ML, робототехника, CAD-моделирование, администрирование, техническое зрение, VR/AR}

\end{document}